\documentclass[border=10pt]{standalone}

% --- packages ---
\usepackage{tikz}
\usetikzlibrary{positioning, arrows.meta, shapes.geometric, fit, backgrounds, calc}

% --- palette -------------------------------------------------------------
% The five OpenTikZ palette names, re-bound to the area-exam deck's part colours
% so this figure matches the other 71 slides. Body code uses names only.
%   otblue  = Part 1   otteal  = Part 2   otorange = Part 3
%   otpurple = auxiliary / "travels back out" / slow loop   otgray = structure
\definecolor{otblue}{HTML}{4472C4}
\definecolor{otteal}{HTML}{1F8A44}
\definecolor{otorange}{HTML}{C8243C}
\definecolor{otpurple}{HTML}{6B4FA8}
\definecolor{otgray}{HTML}{5A5A5A}

% --- shared glyph paths ---------------------------------------------------
\def\otbrainpath{(-0.40,0.02)
  .. controls (-0.46,0.34) and (-0.22,0.50) .. (0,0.44)
  .. controls (0.22,0.50) and (0.46,0.34) .. (0.40,0.02)
  .. controls (0.46,-0.30) and (0.20,-0.46) .. (0,-0.40)
  .. controls (-0.20,-0.46) and (-0.46,-0.30) .. (-0.40,0.02) -- cycle}

% ==== icon pics ===========================================================
\tikzset{
  % --- a brain: outline + midline + two folds -----------------------------
  picbrain/.pic={
    \filldraw[draw=otgray!80, fill=otpurple!12, line width=0.9pt] \otbrainpath;
    \draw[otgray!65, line width=0.7pt] (0,0.44) -- (0,-0.40);
    \draw[otgray!55, line width=0.55pt] (-0.30,0.28) .. controls (-0.12,0.20) and (-0.12,0.10) .. (-0.30,0.06);
    \draw[otgray!55, line width=0.55pt] (-0.32,-0.06) .. controls (-0.14,-0.12) and (-0.14,-0.22) .. (-0.28,-0.28);
    \draw[otgray!55, line width=0.55pt] ( 0.30,0.28) .. controls ( 0.12,0.20) and ( 0.12,0.10) .. ( 0.30,0.06);
    \draw[otgray!55, line width=0.55pt] ( 0.32,-0.06) .. controls ( 0.14,-0.12) and ( 0.14,-0.22) .. ( 0.28,-0.28);
  },
  % --- HyBRiD: a hyperedge (blob) spanning three of five regions ----------
  % two overlapping hyperedges: K parallel heads, each spanning its own subset
  pichyper/.pic={
    \filldraw[draw=otblue!75, fill=otblue!22, line width=1pt, rounded corners=8pt]
      (-0.36,0.30) -- (0.30,0.36) -- (0.02,-0.20) -- cycle;
    \filldraw[draw=otorange!70, fill=otorange!16, line width=1pt, rounded corners=8pt]
      (0.02,-0.20) -- (0.46,-0.06) -- (-0.34,-0.34) -- cycle;
    \foreach \p in {(-0.36,0.30),(0.30,0.36),(0.02,-0.20),(-0.34,-0.34),(0.46,-0.06)}
      \filldraw[draw=otblue!85!black, fill=white, line width=1.1pt] \p circle[radius=0.095];
  },
  % --- Neuro-attn: a brain parcellated by an atlas ------------------------
  picatlas/.pic={
    \begin{scope}
      \clip \otbrainpath;
      \fill[otblue!34] (-0.6,0.02) rectangle (0,0.6);
      \fill[otblue!14] (0,0.02) rectangle (0.6,0.6);
      \fill[otteal!24] (-0.6,-0.6) rectangle (0,0.02);
      \fill[otorange!22] (0,-0.6) rectangle (0.6,0.02);
      \draw[white, line width=0.9pt] (-0.6,0.02) -- (0.6,0.02);
      \draw[white, line width=0.9pt] (0,0.6) -- (0,-0.6);
    \end{scope}
    \draw[otgray!80, line width=0.9pt] \otbrainpath;
    \filldraw[draw=otblue!80!black, fill=white, line width=0.7pt] (-0.22,0.22) circle[radius=0.07];
    \filldraw[draw=otblue!80!black, fill=white, line width=0.7pt] (0.20,-0.18) circle[radius=0.07];
  },
  % --- FA-EEG: two mirror-image traces about a dashed axis ----------------
  picmirror/.pic={
    \draw[otgray!45, dashed, line width=0.6pt] (0,0.46) -- (0,-0.46);
    \draw[otblue!80!black, line width=1pt]
      (-0.46,-0.10) .. controls (-0.36,0.34) and (-0.26,-0.36) .. (-0.16,0.10)
                    .. controls (-0.10,0.32) .. (-0.05,-0.02);
    \draw[otblue!80!black, line width=1pt]
      ( 0.46,-0.10) .. controls ( 0.36,0.34) and ( 0.26,-0.36) .. ( 0.16,0.10)
                    .. controls ( 0.10,0.32) .. ( 0.05,-0.02);
    \draw[otgray!55, line width=0.5pt, -{Stealth[length=1.3mm]}] (-0.30,-0.36) -- (-0.06,-0.36);
    \draw[otgray!55, line width=0.5pt, -{Stealth[length=1.3mm]}] ( 0.30,-0.36) -- ( 0.06,-0.36);
  },
  % --- SWA: two clusters plus one long bridge (small-worldness) ----------
  picsmallworld/.pic={
    \foreach \a/\b in {-0.42/0.22,-0.20/0.34,-0.26/0.02}{
      \foreach \c/\d in {-0.42/0.22,-0.20/0.34,-0.26/0.02}{
        \draw[otpurple!55, line width=0.8pt] (\a,\b) -- (\c,\d);}}
    \foreach \a/\b in {0.22/-0.06,0.44/0.10,0.34/-0.30}{
      \foreach \c/\d in {0.22/-0.06,0.44/0.10,0.34/-0.30}{
        \draw[otpurple!55, line width=0.8pt] (\a,\b) -- (\c,\d);}}
    \draw[otorange!90, line width=1.4pt] (-0.20,0.34) -- (0.22,-0.06);
    \foreach \p in {(-0.42,0.22),(-0.20,0.34),(-0.26,0.02),(0.22,-0.06),(0.44,0.10),(0.34,-0.30)}
      \filldraw[draw=otpurple!80!black, fill=otpurple!28, line width=0.6pt] \p circle[radius=0.075];
  },
  % --- fMRI: a voxel grid with two active cells ---------------------------
  picfmri/.pic={
    \foreach \i in {0,1,2}{\foreach \j in {0,1,2}{
      \draw[otgray!40, line width=0.4pt] (-0.42+0.28*\i,-0.42+0.28*\j) rectangle ++(0.28,0.28);}}
    \fill[otteal!65] (-0.14,-0.14) rectangle ++(0.28,0.28);
    \fill[otteal!32] (0.14,0.14) rectangle ++(0.28,0.28);
    \fill[otteal!22] (-0.42,-0.42) rectangle ++(0.28,0.28);
  },
  % --- PRISM-I: which space is fMRI closer to? (a tilted balance) ---------
  picbalance/.pic={
    \draw[otgray!80, line width=1.3pt] (0,-0.30) -- (0,0.16);
    \draw[otgray!80, line width=1.3pt] (-0.15,-0.34) -- (0.15,-0.34);
    \draw[otgray!80, line width=1.3pt] (-0.46,0.00) -- (0.46,0.26);
    \draw[otgray!70, line width=0.9pt] (-0.46,0.00) -- (-0.46,-0.06);
    \filldraw[draw=otteal!80!black, fill=otteal!30, line width=1.1pt, rounded corners=1pt]
      (-0.62,-0.36) rectangle (-0.30,-0.06);
    \foreach \y in {-0.14,-0.23}{\draw[otteal!85!black, line width=1.1pt] (-0.57,\y) -- (-0.35,\y);}
    \draw[otgray!70, line width=0.9pt] (0.46,0.26) -- (0.46,0.18);
    \filldraw[draw=otgray!70, fill=otgray!12, line width=1.1pt, rounded corners=1pt]
      (0.30,-0.12) rectangle (0.62,0.18);
    \fill[otgray!55] (0.35,-0.06) -- (0.46,0.09) -- (0.57,-0.06) -- cycle;
  },
  % --- a page of text (the target space) ---------------------------------
  pictext/.pic={
    \filldraw[draw=otgray!65, fill=white, line width=0.7pt, rounded corners=2pt]
      (-0.30,-0.42) rectangle (0.30,0.42);
    \foreach \y in {0.26,0.12,-0.02,-0.16}{\draw[otteal!75!black, line width=0.9pt] (-0.20,\y) -- (0.20,\y);}
    \draw[otteal!75!black, line width=0.9pt] (-0.20,-0.30) -- (0.04,-0.30);
  },
  % --- PRISM-II: object-centric reconstruction (objects kept apart) ------
  % solid, not dashed: dash patterns disappear at the size this renders.
  picobjects/.pic={
    \filldraw[draw=otgray!70, fill=white, line width=1pt, rounded corners=2pt]
      (-0.46,-0.36) rectangle (0.46,0.36);
    \filldraw[draw=otteal!80!black, fill=otteal!45, line width=1pt]
      (-0.21,0.04) circle[radius=0.21];
    \filldraw[draw=otorange!80!black, fill=otorange!40, line width=1pt, rounded corners=1pt]
      (0.06,-0.28) rectangle (0.38,0.10);
  },
  % --- D-Edit: one image, one prompt per item ----------------------------
  picitems/.pic={
    \fill[otteal!28] (-0.44,-0.32) rectangle (-0.08,0.32);
    \fill[otorange!30] (-0.04,-0.32) rectangle (0.18,0.06);
    \fill[otpurple!26] (0.22,-0.32) rectangle (0.44,0.32);
    \fill[otgray!10] (-0.04,0.10) rectangle (0.18,0.32);
    \draw[otgray!75, line width=1pt, rounded corners=2pt] (-0.46,-0.34) rectangle (0.46,0.34);
    \foreach \x in {-0.06,0.20}{\draw[white, line width=1.6pt] (\x,-0.34) -- (\x,0.34);}
  },
  % --- Mind World Model: EEG drives the prediction of the next frame -----
  picmindworld/.pic={
    \draw[otorange!85!black, line width=1pt]
      (-0.86,0.00) .. controls (-0.78,0.26) and (-0.70,-0.26) .. (-0.62,0.04)
                   .. controls (-0.54,0.30) and (-0.46,-0.30) .. (-0.38,0.02)
                   .. controls (-0.32,0.20) .. (-0.26,-0.04);
    \draw[otgray!70, line width=0.7pt, -{Stealth[length=1.5mm]}] (-0.20,0.00) -- (-0.02,0.00);
    \filldraw[draw=otgray!60, fill=otgray!10, line width=0.7pt, rounded corners=1.5pt]
      (0.04,-0.22) rectangle (0.38,0.22);
    \filldraw[draw=otorange!80!black, fill=otorange!14, line width=0.8pt, rounded corners=1.5pt]
      (0.46,-0.26) rectangle (0.86,0.26);
    \draw[otgray!70, line width=0.7pt, -{Stealth[length=1.4mm]}] (0.40,0.00) -- (0.44,0.00);
    \fill[otgray!35] (0.09,-0.17) -- (0.18,-0.02) -- (0.27,-0.17) -- cycle;
    \fill[otorange!45] (0.53,-0.20) -- (0.64,-0.01) -- (0.75,-0.20) -- cycle;
  },
  % --- DySta: one static row reused, dynamic row recomputed --------------
  pictokens/.pic={
    \foreach \i in {0,1,2,3}{
      \filldraw[draw=otorange!80!black, fill=otorange!35, line width=0.9pt, rounded corners=1pt]
        (-0.46+0.24*\i,0.10) rectangle ++(0.18,0.18);}
    \draw[otgray!55, dashed, line width=0.8pt] (-0.54,0.02) -- (0.54,0.02);
    \foreach \i in {0,1,2,3}{
      \filldraw[draw=otgray!70, fill=white, line width=0.9pt, rounded corners=1pt]
        (-0.46+0.24*\i,-0.30) rectangle ++(0.18,0.18);}
    \draw[otorange!80!black, line width=0.7pt, -{Stealth[length=1.3mm]}]
      (0.30,0.19) .. controls (0.52,0.19) and (0.52,0.19) .. (0.52,0.10);
  },
  % --- a robot arm (the embodied end) ------------------------------------
  picrobot/.pic={
    \fill[otgray!65] (-0.42,-0.44) rectangle (-0.16,-0.34);
    \draw[otgray!70, line width=1.5pt, line cap=round] (-0.29,-0.34) -- (-0.29,-0.06);
    \draw[otgray!70, line width=1.5pt, line cap=round] (-0.29,-0.06) -- (0.06,0.24);
    \draw[otorange!85!black, line width=1.5pt, line cap=round] (0.06,0.24) -- (0.36,0.06);
    \draw[otorange!85!black, line width=1.1pt, line cap=round] (0.36,0.06) -- (0.48,0.16);
    \draw[otorange!85!black, line width=1.1pt, line cap=round] (0.36,0.06) -- (0.46,-0.06);
    \filldraw[draw=otgray!80, fill=white, line width=1pt] (-0.29,-0.06) circle[radius=0.07];
    \filldraw[draw=otgray!80, fill=white, line width=1pt] (0.06,0.24) circle[radius=0.07];
  },
  % --- PHIL: a sparse goal, a planned path -------------------------------
  picgoal/.pic={
    \draw[otpurple!85, line width=1.3pt] (0.30,0.18) circle[radius=0.22];
    \fill[otpurple!85] (0.30,0.18) circle[radius=0.09];
    \draw[otpurple!85, line width=1.3pt, line cap=round, line join=round]
      (-0.50,-0.28) -- (-0.20,-0.06) -- (-0.02,-0.28) -- (0.30,0.18);
    \filldraw[draw=otpurple!85!black, fill=white, line width=1.1pt] (-0.50,-0.28) circle[radius=0.10];
  },
  % --- the world -----------------------------------------------------------
  picglobe/.pic={
    \filldraw[draw=otteal!80!black, fill=otteal!10, line width=0.9pt] (0,0) circle[radius=0.40];
    \draw[otteal!80!black, line width=0.6pt] (-0.40,0) -- (0.40,0);
    \draw[otteal!80!black, line width=0.6pt] (-0.31,0.25) -- (0.31,0.25);
    \draw[otteal!80!black, line width=0.6pt] (-0.31,-0.25) -- (0.31,-0.25);
    \draw[otteal!80!black, line width=0.6pt] (0,0.40) .. controls (-0.24,0.12) and (-0.24,-0.12) .. (0,-0.40);
    \draw[otteal!80!black, line width=0.6pt] (0,0.40) .. controls ( 0.24,0.12) and ( 0.24,-0.12) .. (0,-0.40);
  },
}

\begin{document}
% ==== parameters (edit these) ============================================
\def\colw{9.4}        % width of one part column (cm)
\def\colgap{2.3}      % gap between part columns (cm)
\def\rowa{0}          % y of row A inside a column
\def\rowb{-2.4}      % y of row B
\def\rowc{-4.75}      % y of row C
\def\rowbb{-3.55}     % y of the two Part 3 branches
\def\bandy{6.9}       % y of the premise band
\def\cardw{2.6}       % width of a small (three-across) card
\def\cardh{2.3}       % height of a row-A card
\def\iconup{0.30}     % icon offset above a card's centre
\def\labdn{-0.48}     % label offset below a card's centre
% labels
\def\partonetitle{PART 1 \textbf{The Foundation}}
\def\parttwotitle{PART 2 \textbf{The Bridge}}
\def\partthreetitle{PART 3 \textbf{The Embodiment}}
\def\partonesub{Models for Brain Signals}
\def\parttwosub{From Signals to Representations}
\def\partthreesub{From Representations to Interactions}
% =========================================================================

\begin{tikzpicture}[
    >={Stealth[length=2.2mm]},
    font=\sffamily,
    card/.style={
      rounded corners=4pt, line width=0.8pt, align=center,
      inner sep=6pt, font=\sffamily\fontsize{12}{14}\selectfont,
    },
    p1/.style={card, draw=otblue!75!black, fill=otblue!10},
    p2/.style={card, draw=otteal!75!black, fill=otteal!10},
    p3/.style={card, draw=otorange!75!black, fill=otorange!08},
    ghost/.style={card, draw=otpurple!70, fill=otpurple!07, dashed},
    bar/.style={card, draw=otgray!45, fill=white, font=\sffamily\fontsize{13}{15}\selectfont},
    hub/.style={card, draw=otblue!60!black, fill=otblue!06, line width=1pt,
                font=\sffamily\fontsize{13}{15}\selectfont, text=otgray!20!black},
    clab/.style={font=\sffamily\fontsize{12}{14}\selectfont\bfseries, text=otgray!25!black, align=center},
    csub/.style={font=\sffamily\fontsize{10.5}{12}\selectfont, text=otgray!85, align=center},
    link/.style={draw=otgray!75, line width=1pt, ->},
    backlink/.style={draw=otpurple!85, line width=1pt, ->, dashed},
    elab/.style={font=\sffamily\fontsize{10.5}{12}\selectfont, text=otgray!95, inner sep=2pt, midway, right, xshift=3pt},
    sub/.style={font=\sffamily\fontsize{11}{13}\selectfont, text=otgray!85},
    group/.style={draw=otgray!35, rounded corners=6pt, line width=0.8pt,
                  inner sep=14pt, fill=otgray!4},
  ]

  % ---------- column origins ----------
  \def\cA{0}
  \pgfmathsetmacro{\cB}{\colw+\colgap}
  \pgfmathsetmacro{\cC}{2*(\colw+\colgap)}

  % ================= PART 1 : discover -> build in -> travel back out =====
  \node[p1, minimum width=\cardw cm, minimum height=\cardh cm] (hybrid) at (\cA-2.9,\rowa) {};
    \pic[scale=1.05] at ([yshift=\iconup cm]hybrid.center) {pichyper};
    \node[clab] at ([yshift=\labdn cm]hybrid.center) {HyBRiD};

  \node[p1, minimum width=\cardw cm, minimum height=\cardh cm] (neuro) at (\cA,\rowa) {};
    \pic[scale=1.05] at ([yshift=\iconup cm]neuro.center) {picatlas};
    \node[clab] at ([yshift=\labdn cm]neuro.center) {Neuro-attn};

  \node[p1, minimum width=\cardw cm, minimum height=\cardh cm] (faeeg) at (\cA+2.9,\rowa) {};
    \pic[scale=1.05] at ([yshift=\iconup cm]faeeg.center) {picmirror};
    \node[clab] at ([yshift=\labdn cm]faeeg.center) {FA-EEG};

  % the hub: the prior itself. HyBRiD writes into it, the other two read from it.
  \node[hub, minimum width=8.4cm, minimum height=1cm] (foundation) at (\cA,\rowb)
    {\textbf{neural priors} --- how the brain is organised};

  \node[ghost, minimum width=8.4cm, minimum height=1.4cm] (backout) at (\cA,\rowc) {};
    \pic[scale=1.1] at ([xshift=-3.2cm]backout.center) {picsmallworld};
    \node[clab, anchor=west, text=otpurple!35!black] at ([xshift=-2.5cm,yshift=0cm]backout.center) {SWA};

  % HyBRiD discovers a prior; the other two have one built in. Opposite directions.
  \draw[link] (hybrid.south) -- ([xshift=-2.9cm]foundation.north)
    node[elab]{ML $\to$ prior};
  \draw[link] (foundation.north) -- (neuro.south)
    node[elab]{prior $\to$ ML};
  \draw[link] ([xshift=2.9cm]foundation.north) -- (faeeg.south);
  \draw[backlink] (foundation.south) -- (backout.north)
    node[elab, text=otpurple]{the prior travels back out};

  % ================= PART 2 : four linked research questions =============
  % The arrows express motivation, not an integrated set of model weights.
  \node[p2, minimum width=4.05cm, minimum height=\cardh cm] (prism1) at (\cB-2.18,\rowa) {};
    \pic[scale=1.0] at ([yshift=0.42cm]prism1.center) {picbalance};
    \node[clab] at ([yshift=-0.46cm]prism1.center) {PRISM I};
    \node[csub] at ([yshift=-0.84cm]prism1.center) {Choose alignment space};

  \node[p2, minimum width=4.05cm, minimum height=\cardh cm] (encoder) at (\cB+2.18,\rowa) {};
    \pic[scale=0.95] at ([yshift=0.42cm]encoder.center) {picfmri};
    \node[clab] at ([yshift=-0.46cm]encoder.center) {MindLLM};
    \node[csub] at ([yshift=-0.84cm]encoder.center) {Align fMRI to language};
  \draw[link] (prism1.east) -- (encoder.west);

  \node[p2, minimum width=8.4cm, minimum height=1.6cm] (space) at (\cB,\rowb) {};
    \pic[scale=1.0] at ([xshift=-3.3cm]space.center) {pictext};
    \node[anchor=west, font=\sffamily\fontsize{12}{14}\selectfont, align=left, text=otgray!20!black]
      at ([xshift=-2.5cm]space.center)
      {\textbf{Language representation}\\Read out multiple brain functions};
  \draw[link] (encoder.south) -- ([xshift=2.18cm]space.north);

  \node[p2, minimum width=4.05cm, minimum height=1.8cm] (prism2) at (\cB-2.18,\rowc) {};
    \node[clab] at ([yshift=0.38cm]prism2.center) {PRISM II};
    \node[csub, font=\sffamily\fontsize{9}{10.5}\selectfont] at ([yshift=-0.25cm]prism2.center) {Language-to-vision\\Object-centric generation};
  \node[p2, minimum width=4.05cm, minimum height=1.8cm] (dedit) at (\cB+2.18,\rowc) {};
    \node[clab] at ([yshift=0.38cm]dedit.center) {D-Edit};
    \node[csub, font=\sffamily\fontsize{9}{10.5}\selectfont] at ([yshift=-0.25cm]dedit.center) {Item prompts + attention\\Stronger disentanglement};
  \draw[link] ([xshift=-2.18cm]space.south) -- (prism2.north);
  \draw[link] (prism2.east) -- (dedit.west);

  % ================= PART 3 : MWM parallel to (PHIL parallel to DySta) ====
  % Nested groups express research organization, not model input/output.
  \node[p3, minimum width=4.05cm, minimum height=5.8cm] (mwm) at (\cC-2.18,-1.75) {};
    \node[clab] at ([yshift=2.14cm]mwm.center) {MWM};
    \pic[scale=1.6] at ([yshift=0.68cm]mwm.center) {picmindworld};
    \node[csub, font=\sffamily\fontsize{9}{10.5}\selectfont] at ([yshift=-1.1cm]mwm.center) {Brain signals\\to action meaning\\[4pt]EEG latent actions\\predict visual consequences};

  \node[p3, minimum width=4.05cm, minimum height=5.8cm, fill=white] (goalbranch) at (\cC+2.18,-1.75) {};
    \node[clab] at ([yshift=2.14cm]goalbranch.center) {Act toward\\specified goals};
  \node[ghost, minimum width=3.45cm, minimum height=1.65cm] (phil) at (\cC+2.18,-1.45) {};
    \node[clab, text=otpurple!35!black] at ([yshift=0.34cm]phil.center) {PHIL};
    \node[csub, font=\sffamily\fontsize{9}{10.5}\selectfont] at ([yshift=-0.26cm]phil.center) {Learned heuristic\\for graph search};
  \node[p3, minimum width=3.45cm, minimum height=1.65cm] (dysta) at (\cC+2.18,-3.5) {};
    \node[clab] at ([yshift=0.34cm]dysta.center) {DySta};
    \node[csub, font=\sffamily\fontsize{9}{10.5}\selectfont] at ([yshift=-0.26cm]dysta.center) {Visual policy\\Memory + fast inference};
  \node[csub] (p3scope) at (\cC,-5.05) {Separately evaluated capabilities};

  % ================= part boundaries + headers ============================
  \begin{scope}[on background layer]
    \node[group, fit=(hybrid)(neuro)(faeeg)(foundation)(backout)]      (gA) {};
    \node[group, fit=(prism1)(encoder)(space)(prism2)(dedit)]          (gB) {};
    \node[group, fit=(mwm)(goalbranch)(p3scope)]                              (gC) {};
  \end{scope}
  \foreach \g/\c/\t/\s in {%
      gA/otblue/{\partonetitle}/{\partonesub},
      gB/otteal/{\parttwotitle}/{\parttwosub},
      gC/otorange/{\partthreetitle}/{\partthreesub}}
    {%
      \node[fill=\c, text=white, inner xsep=11pt, inner ysep=6pt, align=left,
            anchor=south west, font=\sffamily\fontsize{13.5}{15.5}\selectfont, rounded corners=3pt]
        (\g h) at ([yshift=4pt]\g.north west)
        {\t\\[1.5pt]\textcolor{white!88}{\fontsize{10.5}{12}\selectfont\s}};
    }
  \draw[link, line width=1.3pt] (gA.east) -- (gB.west)
    node[midway, above, font=\sffamily\fontsize{10.5}{12}\selectfont, text=otgray, align=center]{signals\\$\rightarrow$ semantics};
  \draw[link, line width=1.3pt] (gB.east) -- (gC.west)
    node[midway, above, font=\sffamily\fontsize{10.5}{12}\selectfont, text=otgray, align=center]{semantics\\$\rightarrow$ action};

  % ================= premise band: a lossy channel ========================
  \coordinate (bC) at (\cB,\bandy);
  \coordinate (f1) at ($(bC)+(-5.0,0)$);
  \coordinate (f2) at ($(bC)+(-1.3,0)$);
  \coordinate (f3) at ($(bC)+(1.3,0)$);
  \coordinate (f4) at ($(bC)+(5.0,0)$);

  % the pipe: wide in, pinched at language, narrower out
  \filldraw[draw=otgray!45, fill=otgray!8, line width=0.9pt]
    ($(f1)+(0,1.05)$) -- ($(f2)+(0,0.26)$) -- ($(f3)+(0,0.26)$) -- ($(f4)+(0,0.52)$) --
    ($(f4)+(0,-0.52)$) -- ($(f3)+(0,-0.26)$) -- ($(f2)+(0,-0.26)$) -- ($(f1)+(0,-1.05)$) -- cycle;
  \draw[otgray!55, line width=1.2pt, ->] ($(f1)+(0.2,0)$) -- ($(f4)+(-0.25,0)$);
  \filldraw[draw=otorange, fill=white, line width=1.3pt, rounded corners=3pt]
    ($(f2)+(-0.12,-0.58)$) rectangle ($(f3)+(0.12,0.58)$);
  \node[font=\sffamily\fontsize{17}{19}\selectfont\bfseries, text=otorange] at (bC) {language};

  % information that does not survive the pinch
  \foreach \dx/\dy/\s in {-0.90/-1.02/0.20, -0.42/-1.34/0.26, 0.08/-0.98/0.30,
                          0.54/-1.44/0.22, 1.00/-1.06/0.25}{
    \filldraw[draw=otorange!55, fill=otorange!22, line width=0.6pt, rounded corners=1pt]
      ($(bC)+(\dx,\dy)$) rectangle ++(\s,\s);}
  \node[font=\sffamily\fontsize{10.5}{12}\selectfont\itshape, text=otorange!55!black, anchor=north]
    at ($(bC)+(0.35,-1.80)$) {what the words leave out};

  \pic[scale=1.7] at ($(f1)+(-1.5,0)$) {picbrain};
  \node[font=\sffamily\fontsize{21}{23}\selectfont\bfseries, text=otgray!15!black] at ($(f1)+(-1.5,-1.35)$) {MIND};
  \pic[scale=1.7] at ($(f4)+(1.5,0)$) {picglobe};
  \node[font=\sffamily\fontsize{21}{23}\selectfont\bfseries, text=otgray!15!black] at ($(f4)+(1.5,-1.35)$) {WORLD};

  \node[sub, anchor=north] at ($(bC)+(0,-2.50)$)
    {the channel is lossy in both directions --- measured in RepQA};
  \node[font=\sffamily\fontsize{13}{15}\selectfont\bfseries, text=otgray!15!black, anchor=north]
    at ($(bC)+(0,-2.96)$)
    {So: connect brain signals without requiring verbal instructions};

  % fan out from the premise into the three parts
  \coordinate (hub) at ($(bC)+(0,-3.52)$);
  \foreach \g in {gA,gB,gC}{
    \draw[otgray!45, line width=0.9pt, -{Stealth[length=2mm]}, rounded corners=8pt]
      (hub) -| ($(\g.north)+(2.6,0.16)$);}

  % ================= closing bar: the bypass the three parts build =========
  \coordinate (byL) at ($(gA.south west)+(0.4,-0.75)$);
  \coordinate (byR) at ($(gC.south east)+(-0.4,-0.75)$);
  \draw[-{Stealth[length=5mm,width=4mm]}, line width=3.5pt, otgray!30] (byL) -- (byR);
  \node[font=\sffamily\fontsize{13}{15}\selectfont\bfseries, text=otgray!15!black, fill=white, inner sep=5pt]
    at ($(byL)!0.5!(byR)$)
    {brain signals, semantics, and behavior --- complementary research capabilities};

\end{tikzpicture}
\end{document}
